% LatentShift: Theoretical Analysis
% This file contains the formal propositions and proofs for the paper.
% Include via % LatentShift: Theoretical Analysis
% This file contains the formal propositions and proofs for the paper.
% Include via % LatentShift: Theoretical Analysis
% This file contains the formal propositions and proofs for the paper.
% Include via % LatentShift: Theoretical Analysis
% This file contains the formal propositions and proofs for the paper.
% Include via \input{theory} from the main paper file.

\section{Theoretical Analysis}
\label{sec:theory}

We present formal guarantees for the LatentShift method. All proofs follow
from standard linear algebra and the properties of orthogonal projections.

\subsection{Setup and Notation}

Let $f_\theta: \mathcal{X} \to \mathbb{R}^d$ denote the encoder mapping
inputs to a $d$-dimensional latent space. After training on task $t$, we
collect latent activations $Z_t = \{f_\theta(x) : x \in \mathcal{D}_t\}$
and compute via SVD a rank-$r_t$ orthonormal basis $V_t \in \mathbb{R}^{d \times r_t}$
capturing at least $\tau$ fraction of the variance (where $\tau$ is the
threshold hyperparameter, typically 0.99).

The \emph{archive basis} after $T$ tasks is the orthonormalized union
$A_T \in \mathbb{R}^{d \times k_T}$, where $k_T \leq \sum_{t=1}^T r_t$,
obtained by QR decomposition of $[A_{T-1} \mid V_T]$.

The \emph{free-subspace projector} is:
\begin{equation}
    P = I_d - A_T A_T^\top
    \label{eq:projector}
\end{equation}

During training on task $T+1$, all gradients $g$ flowing through the latent
representation are projected as $g' = Pg$.

\subsection{Proposition 1: Zero-Forgetting Guarantee}

\begin{proposition}[Zero Forgetting]
\label{prop:zero-forgetting}
Let $A \in \mathbb{R}^{d \times k}$ be the archive basis with orthonormal
columns ($A^\top A = I_k$), and let $P = I_d - AA^\top$ be the free-subspace
projector. For any archived representation $z \in \mathrm{col}(A)$ and any
projected gradient $g' = Pg$:
\begin{equation}
    z^\top g' = 0
\end{equation}
That is, the projected gradient has zero component along any direction in
the archive subspace. Consequently, parameter updates driven by $g'$
cannot modify the encoder's output for any input whose latent
representation lies in $\mathrm{col}(A)$.
\end{proposition}

\begin{proof}
Since $z \in \mathrm{col}(A)$, there exists $\alpha \in \mathbb{R}^k$ such
that $z = A\alpha$. Then:
\begin{align}
    z^\top g' &= (A\alpha)^\top (I - AA^\top) g \\
              &= \alpha^\top A^\top g - \alpha^\top \underbrace{A^\top A}_{= I_k} A^\top g \\
              &= \alpha^\top A^\top g - \alpha^\top A^\top g \\
              &= 0
\end{align}
The first-order Taylor expansion of the encoder output change is
$\Delta f_\theta(x) \approx J_z \cdot g'$, where $J_z$ is the Jacobian.
Since $g'$ lies entirely in $\mathrm{null}(A^\top)$, and the archived
representations span $\mathrm{col}(A)$, the inner product
$z^\top \Delta f_\theta(x) = 0$ to first order.
\end{proof}

\begin{remark}[First-Order Approximation]
\label{rem:first-order}
Proposition~\ref{prop:zero-forgetting} guarantees \emph{exact} orthogonality
$z^\top g' = 0$ at the gradient level, but the claim that archived
representations are unchanged relies on a first-order Taylor expansion of
the encoder $f_\theta$. When $f_\theta$ is a deep nonlinear network,
higher-order terms in the parameter update $\Delta \theta$ can cause
small drifts in the latent representation $z = f_\theta(x)$ even though
the projected gradient $g'$ is orthogonal to the archive subspace.
We quantify this residual in Proposition~\ref{prop:second-order}.
\end{remark}

\begin{remark}[Latent-to-Parameter Gap]
\label{rem:latent-to-parameter}
Proposition~\ref{prop:zero-forgetting} establishes orthogonality in the
latent space: $z^\top g' = 0$. To connect this to parameter space, note
that the encoder $f_\theta$ maps parameters $\theta$ to latent
representations $z = f_\theta(x)$. Let $J_\theta = \partial f_\theta(x) /
\partial \theta$ denote the Jacobian of the encoder with respect to
parameters (note: this is distinct from $J_z$ used in the proof above,
which informally denotes the Jacobian of the encoder output change with
respect to latent-space perturbations). If $g' = Pg$ is the projected
gradient in latent space, then by the chain rule, the effective parameter
gradient becomes:
\begin{equation}
    \frac{\partial \mathcal{L}}{\partial \theta}\bigg|_{\text{projected}}
    = J_\theta^\top \, g'
    = J_\theta^\top \, P \, g
    \label{eq:latent-to-param}
\end{equation}
Since $g' = Pg \in \mathrm{null}(A^\top)$ by construction, the parameter
update $\Delta \theta = -\eta \, J_\theta^\top g'$ inherits the
orthogonality constraint: for any archived representation
$z = A\alpha \in \mathrm{col}(A)$, the first-order change in $z$ is
$\Delta z \approx J_\theta \, \Delta \theta = -\eta \, J_\theta J_\theta^\top g'$.
The component of $\Delta z$ along the archive is
$A^\top \Delta z = -\eta \, A^\top J_\theta J_\theta^\top g'$, which
vanishes to first order when $g' \in \mathrm{null}(A^\top)$ and
$J_\theta J_\theta^\top$ preserves this null-space structure (as holds
exactly when $J_\theta$ commutes with $P$, and approximately otherwise).
This is a \emph{first-order} guarantee; the second-order remainder is
captured by Proposition~\ref{prop:second-order}, which bounds the
cumulative drift from nonlinear effects.
\end{remark}

\subsection{Proposition 2: Capacity Bound}

\begin{proposition}[Capacity Bound]
\label{prop:capacity}
For a latent space of dimension $d$, with task $t$ occupying a subspace
of rank $r_t$, the LatentShift method guarantees zero forgetting for all
tasks if and only if:
\begin{equation}
    k_T = \mathrm{rank}(A_T) \leq d
\end{equation}
After QR reorthogonalization, $k_T \leq \sum_{t=1}^T r_t$, with equality
when task subspaces are mutually orthogonal. The maximum number of tasks
with guaranteed zero forgetting is:
\begin{equation}
    T_{\max} = \left\lfloor \frac{d}{\bar{r}} \right\rfloor
\end{equation}
where $\bar{r} = \frac{1}{T}\sum_{t=1}^T r_t$ is the average per-task rank.
When $k_T = d$, the free subspace is empty ($P = 0$) and no further
learning is possible without interference.
\end{proposition}

\begin{proof}
The archive basis $A_T$ has $k_T$ orthonormal columns in $\mathbb{R}^d$.
The free projector $P = I - A_T A_T^\top$ has rank $d - k_T$. For $P$
to be non-trivial (allowing new learning), we need $d - k_T > 0$.

Each task $t$ contributes at most $r_t$ new linearly independent directions.
The QR step ensures orthonormality and removes redundancy, so
$k_T \leq \min(\sum_t r_t, d)$. The bound is tight when task subspaces
are orthogonal.

Setting $k_T = d$ yields $P = 0$, meaning $g' = 0$ for all gradients,
preventing any parameter updates. Thus $T_{\max} = \lfloor d/\bar{r} \rfloor$.
\end{proof}

\subsection{Proposition 3: Isometry Preservation}

\begin{proposition}[Free-Subspace Isometry]
\label{prop:isometry}
The projector $P = I - AA^\top$ preserves inner products within the free
subspace. That is, for any $u, v \in \mathrm{null}(A^\top)$:
\begin{equation}
    \langle Pu, Pv \rangle = \langle u, v \rangle
\end{equation}
Moreover, $P$ preserves norms: $\|Pu\| = \|u\|$ for $u \in \mathrm{null}(A^\top)$.
\end{proposition}

\begin{proof}
If $u \in \mathrm{null}(A^\top)$, then $A^\top u = 0$, so
$Pu = u - AA^\top u = u - A \cdot 0 = u$. Similarly $Pv = v$. Therefore:
\[
    \langle Pu, Pv \rangle = \langle u, v \rangle.
\]
For norms: $\|Pu\|^2 = \langle Pu, Pu \rangle = \langle u, u \rangle = \|u\|^2$.
\end{proof}

This ensures that LatentShift does not distort the geometry of new task
representations — learning in the free subspace proceeds as if the
archive dimensions simply did not exist.

\subsection{Proposition 4: Lossy Compression Bound}

\begin{proposition}[Compression Error Bound]
\label{prop:compression}
Suppose the archive has rank $k$ with basis $A = [a_1 \mid \cdots \mid a_k]$
ordered by decreasing importance (singular values $\sigma_1 \geq \cdots
\geq \sigma_k$ from the cumulative SVD). Compressing to rank $k' < k$ by
retaining only the first $k'$ columns yields a representation error bounded by:
\begin{equation}
    \max_{z \in \mathrm{col}(A), \|z\|=1} \|z - \hat{z}\| \leq
    \sqrt{\sum_{i=k'+1}^{k} \sigma_i^2 \Big/ \sum_{i=1}^{k} \sigma_i^2}
\end{equation}
where $\hat{z}$ is the projection of $z$ onto the compressed basis
$A' = [a_1 \mid \cdots \mid a_{k'}]$.
\end{proposition}

\begin{proof}
For any unit vector $z \in \mathrm{col}(A)$, write $z = A\alpha$ with
$\|\alpha\| = 1$ (since $A$ has orthonormal columns). The projection onto
$A'$ gives $\hat{z} = A' A'^\top z = A' \alpha_{1:k'}$ where $\alpha_{1:k'}$
are the first $k'$ components. The error is:
\[
    \|z - \hat{z}\|^2 = \|\alpha_{k'+1:k}\|^2 \leq 1
\]
When the archive was built from SVD with singular values $\sigma_i$, the
importance of each direction is proportional to $\sigma_i^2$. Restricting
to the subspace captured by the cumulative SVD, the relative error for
the worst-case archived representation is bounded by the fraction of
discarded variance.
\end{proof}

\subsection{Proposition 5: Second-Order Forgetting Bound}

\begin{proposition}[Cumulative Second-Order Forgetting]
\label{prop:second-order}
Let $f_\theta$ be an $L$-smooth encoder (i.e.\ its Jacobian $J_\theta f$
is Lipschitz with constant $L$). Suppose training on task $T+1$ takes
$N$ gradient steps with learning rate $\eta$, each producing a projected
gradient $g'_n = P g_n$ with $\|g_n\| \leq G$. Then for any input $x$
whose initial representation $z = f_{\theta_0}(x) \in \mathrm{col}(A)$,
the drift after $N$ steps satisfies:
\begin{equation}
    \|f_{\theta_N}(x) - f_{\theta_0}(x)\| \leq
    \underbrace{0}_{\text{first-order}} +
    \frac{L \eta^2 G^2 N(N-1)}{2}
    \label{eq:second-order-bound}
\end{equation}
\end{proposition}

\begin{proof}
By Taylor expansion of the encoder around parameters $\theta_n$:
\[
    f_{\theta_{n+1}}(x) - f_{\theta_n}(x)
    = J_{\theta_n} f(x) \cdot (-\eta g'_n) + R_n
\]
where $\|R_n\| \leq \frac{L}{2} \|\eta g'_n\|^2 \leq \frac{L\eta^2 G^2}{2}$.

Proposition~\ref{prop:zero-forgetting} shows the first-order term
$J_{\theta_n} f(x) \cdot (-\eta g'_n)$ has zero component in
$\mathrm{col}(A)$ at step $n=0$. At subsequent steps the archived
representation may have drifted by the accumulated remainder, giving
a first-order contribution of $O(\eta^2)$ per step. Summing over $N$
steps:
\[
    \|f_{\theta_N}(x) - f_{\theta_0}(x)\|
    \leq \sum_{n=0}^{N-1} \|R_n\| + O(\eta^2 N)
    \leq \frac{L \eta^2 G^2 N(N-1)}{2}
\]
\end{proof}

\begin{corollary}[Small Learning Rate Regime]
\label{cor:small-lr}
As $\eta \to 0$ with fixed total parameter displacement $\eta N = C$
(i.e.\ more steps at smaller rate), the forgetting bound becomes
$O(\eta) \to 0$. In the continuous gradient-flow limit ($\eta \to 0,
N \to \infty, \eta N = C$), the LatentShift method achieves exactly
zero forgetting for archived representations.
\end{corollary}

This explains the empirical gap between the theoretical zero-forgetting
guarantee and the small but non-zero forgetting observed on CIFAR benchmarks:
large learning rates and deep nonlinear encoders introduce non-negligible
second-order terms. The gap shrinks with smaller learning rates or
shallower encoders (as confirmed on Split-MNIST, where forgetting is
$< 1\%$).

\subsection{Connection to Hilbert's Hotel}

The LatentShift method draws an analogy with Hilbert's Grand Hotel paradox:

\begin{itemize}
    \item \textbf{Rooms} correspond to latent dimensions $\{1, 2, \ldots, d\}$.
    \item \textbf{Guests} correspond to task representations (subspace directions).
    \item The \textbf{shift operator} $n \to 2n$ in Hilbert's Hotel corresponds to
          our SVD-based subspace identification followed by orthogonal projection,
          which ``shifts'' old task representations into a protected archive and
          frees orthogonal directions for new guests.
\end{itemize}

Unlike Hilbert's Hotel (which has countably infinite rooms), the latent space
has finite dimension $d$, leading to the capacity bound in Proposition~\ref{prop:capacity}.
This boundedness motivates the lossy compression mechanism
(Proposition~\ref{prop:compression}) as a practical extension.

\textbf{Key distinction from GPM:} While GPM~\cite{saha2021gradient} also
uses orthogonal projection to prevent forgetting, it operates \emph{per-layer},
maintaining separate subspace bases for each network layer. LatentShift operates
in the \emph{holistic latent representation space}, which has several advantages:
(1)~a single, interpretable capacity measure (archive rank vs.\ latent dimension);
(2)~architectural flexibility (the projector is agnostic to encoder structure);
(3)~a cleaner theoretical framework (Propositions~\ref{prop:zero-forgetting}--\ref{prop:compression}
    apply to the unified latent space rather than fragmented per-layer spaces).
 from the main paper file.

\section{Theoretical Analysis}
\label{sec:theory}

We present formal guarantees for the LatentShift method. All proofs follow
from standard linear algebra and the properties of orthogonal projections.

\subsection{Setup and Notation}

Let $f_\theta: \mathcal{X} \to \mathbb{R}^d$ denote the encoder mapping
inputs to a $d$-dimensional latent space. After training on task $t$, we
collect latent activations $Z_t = \{f_\theta(x) : x \in \mathcal{D}_t\}$
and compute via SVD a rank-$r_t$ orthonormal basis $V_t \in \mathbb{R}^{d \times r_t}$
capturing at least $\tau$ fraction of the variance (where $\tau$ is the
threshold hyperparameter, typically 0.99).

The \emph{archive basis} after $T$ tasks is the orthonormalized union
$A_T \in \mathbb{R}^{d \times k_T}$, where $k_T \leq \sum_{t=1}^T r_t$,
obtained by QR decomposition of $[A_{T-1} \mid V_T]$.

The \emph{free-subspace projector} is:
\begin{equation}
    P = I_d - A_T A_T^\top
    \label{eq:projector}
\end{equation}

During training on task $T+1$, all gradients $g$ flowing through the latent
representation are projected as $g' = Pg$.

\subsection{Proposition 1: Zero-Forgetting Guarantee}

\begin{proposition}[Zero Forgetting]
\label{prop:zero-forgetting}
Let $A \in \mathbb{R}^{d \times k}$ be the archive basis with orthonormal
columns ($A^\top A = I_k$), and let $P = I_d - AA^\top$ be the free-subspace
projector. For any archived representation $z \in \mathrm{col}(A)$ and any
projected gradient $g' = Pg$:
\begin{equation}
    z^\top g' = 0
\end{equation}
That is, the projected gradient has zero component along any direction in
the archive subspace. Consequently, parameter updates driven by $g'$
cannot modify the encoder's output for any input whose latent
representation lies in $\mathrm{col}(A)$.
\end{proposition}

\begin{proof}
Since $z \in \mathrm{col}(A)$, there exists $\alpha \in \mathbb{R}^k$ such
that $z = A\alpha$. Then:
\begin{align}
    z^\top g' &= (A\alpha)^\top (I - AA^\top) g \\
              &= \alpha^\top A^\top g - \alpha^\top \underbrace{A^\top A}_{= I_k} A^\top g \\
              &= \alpha^\top A^\top g - \alpha^\top A^\top g \\
              &= 0
\end{align}
The first-order Taylor expansion of the encoder output change is
$\Delta f_\theta(x) \approx J_z \cdot g'$, where $J_z$ is the Jacobian.
Since $g'$ lies entirely in $\mathrm{null}(A^\top)$, and the archived
representations span $\mathrm{col}(A)$, the inner product
$z^\top \Delta f_\theta(x) = 0$ to first order.
\end{proof}

\begin{remark}[First-Order Approximation]
\label{rem:first-order}
Proposition~\ref{prop:zero-forgetting} guarantees \emph{exact} orthogonality
$z^\top g' = 0$ at the gradient level, but the claim that archived
representations are unchanged relies on a first-order Taylor expansion of
the encoder $f_\theta$. When $f_\theta$ is a deep nonlinear network,
higher-order terms in the parameter update $\Delta \theta$ can cause
small drifts in the latent representation $z = f_\theta(x)$ even though
the projected gradient $g'$ is orthogonal to the archive subspace.
We quantify this residual in Proposition~\ref{prop:second-order}.
\end{remark}

\begin{remark}[Latent-to-Parameter Gap]
\label{rem:latent-to-parameter}
Proposition~\ref{prop:zero-forgetting} establishes orthogonality in the
latent space: $z^\top g' = 0$. To connect this to parameter space, note
that the encoder $f_\theta$ maps parameters $\theta$ to latent
representations $z = f_\theta(x)$. Let $J_\theta = \partial f_\theta(x) /
\partial \theta$ denote the Jacobian of the encoder with respect to
parameters (note: this is distinct from $J_z$ used in the proof above,
which informally denotes the Jacobian of the encoder output change with
respect to latent-space perturbations). If $g' = Pg$ is the projected
gradient in latent space, then by the chain rule, the effective parameter
gradient becomes:
\begin{equation}
    \frac{\partial \mathcal{L}}{\partial \theta}\bigg|_{\text{projected}}
    = J_\theta^\top \, g'
    = J_\theta^\top \, P \, g
    \label{eq:latent-to-param}
\end{equation}
Since $g' = Pg \in \mathrm{null}(A^\top)$ by construction, the parameter
update $\Delta \theta = -\eta \, J_\theta^\top g'$ inherits the
orthogonality constraint: for any archived representation
$z = A\alpha \in \mathrm{col}(A)$, the first-order change in $z$ is
$\Delta z \approx J_\theta \, \Delta \theta = -\eta \, J_\theta J_\theta^\top g'$.
The component of $\Delta z$ along the archive is
$A^\top \Delta z = -\eta \, A^\top J_\theta J_\theta^\top g'$, which
vanishes to first order when $g' \in \mathrm{null}(A^\top)$ and
$J_\theta J_\theta^\top$ preserves this null-space structure (as holds
exactly when $J_\theta$ commutes with $P$, and approximately otherwise).
This is a \emph{first-order} guarantee; the second-order remainder is
captured by Proposition~\ref{prop:second-order}, which bounds the
cumulative drift from nonlinear effects.
\end{remark}

\subsection{Proposition 2: Capacity Bound}

\begin{proposition}[Capacity Bound]
\label{prop:capacity}
For a latent space of dimension $d$, with task $t$ occupying a subspace
of rank $r_t$, the LatentShift method guarantees zero forgetting for all
tasks if and only if:
\begin{equation}
    k_T = \mathrm{rank}(A_T) \leq d
\end{equation}
After QR reorthogonalization, $k_T \leq \sum_{t=1}^T r_t$, with equality
when task subspaces are mutually orthogonal. The maximum number of tasks
with guaranteed zero forgetting is:
\begin{equation}
    T_{\max} = \left\lfloor \frac{d}{\bar{r}} \right\rfloor
\end{equation}
where $\bar{r} = \frac{1}{T}\sum_{t=1}^T r_t$ is the average per-task rank.
When $k_T = d$, the free subspace is empty ($P = 0$) and no further
learning is possible without interference.
\end{proposition}

\begin{proof}
The archive basis $A_T$ has $k_T$ orthonormal columns in $\mathbb{R}^d$.
The free projector $P = I - A_T A_T^\top$ has rank $d - k_T$. For $P$
to be non-trivial (allowing new learning), we need $d - k_T > 0$.

Each task $t$ contributes at most $r_t$ new linearly independent directions.
The QR step ensures orthonormality and removes redundancy, so
$k_T \leq \min(\sum_t r_t, d)$. The bound is tight when task subspaces
are orthogonal.

Setting $k_T = d$ yields $P = 0$, meaning $g' = 0$ for all gradients,
preventing any parameter updates. Thus $T_{\max} = \lfloor d/\bar{r} \rfloor$.
\end{proof}

\subsection{Proposition 3: Isometry Preservation}

\begin{proposition}[Free-Subspace Isometry]
\label{prop:isometry}
The projector $P = I - AA^\top$ preserves inner products within the free
subspace. That is, for any $u, v \in \mathrm{null}(A^\top)$:
\begin{equation}
    \langle Pu, Pv \rangle = \langle u, v \rangle
\end{equation}
Moreover, $P$ preserves norms: $\|Pu\| = \|u\|$ for $u \in \mathrm{null}(A^\top)$.
\end{proposition}

\begin{proof}
If $u \in \mathrm{null}(A^\top)$, then $A^\top u = 0$, so
$Pu = u - AA^\top u = u - A \cdot 0 = u$. Similarly $Pv = v$. Therefore:
\[
    \langle Pu, Pv \rangle = \langle u, v \rangle.
\]
For norms: $\|Pu\|^2 = \langle Pu, Pu \rangle = \langle u, u \rangle = \|u\|^2$.
\end{proof}

This ensures that LatentShift does not distort the geometry of new task
representations — learning in the free subspace proceeds as if the
archive dimensions simply did not exist.

\subsection{Proposition 4: Lossy Compression Bound}

\begin{proposition}[Compression Error Bound]
\label{prop:compression}
Suppose the archive has rank $k$ with basis $A = [a_1 \mid \cdots \mid a_k]$
ordered by decreasing importance (singular values $\sigma_1 \geq \cdots
\geq \sigma_k$ from the cumulative SVD). Compressing to rank $k' < k$ by
retaining only the first $k'$ columns yields a representation error bounded by:
\begin{equation}
    \max_{z \in \mathrm{col}(A), \|z\|=1} \|z - \hat{z}\| \leq
    \sqrt{\sum_{i=k'+1}^{k} \sigma_i^2 \Big/ \sum_{i=1}^{k} \sigma_i^2}
\end{equation}
where $\hat{z}$ is the projection of $z$ onto the compressed basis
$A' = [a_1 \mid \cdots \mid a_{k'}]$.
\end{proposition}

\begin{proof}
For any unit vector $z \in \mathrm{col}(A)$, write $z = A\alpha$ with
$\|\alpha\| = 1$ (since $A$ has orthonormal columns). The projection onto
$A'$ gives $\hat{z} = A' A'^\top z = A' \alpha_{1:k'}$ where $\alpha_{1:k'}$
are the first $k'$ components. The error is:
\[
    \|z - \hat{z}\|^2 = \|\alpha_{k'+1:k}\|^2 \leq 1
\]
When the archive was built from SVD with singular values $\sigma_i$, the
importance of each direction is proportional to $\sigma_i^2$. Restricting
to the subspace captured by the cumulative SVD, the relative error for
the worst-case archived representation is bounded by the fraction of
discarded variance.
\end{proof}

\subsection{Proposition 5: Second-Order Forgetting Bound}

\begin{proposition}[Cumulative Second-Order Forgetting]
\label{prop:second-order}
Let $f_\theta$ be an $L$-smooth encoder (i.e.\ its Jacobian $J_\theta f$
is Lipschitz with constant $L$). Suppose training on task $T+1$ takes
$N$ gradient steps with learning rate $\eta$, each producing a projected
gradient $g'_n = P g_n$ with $\|g_n\| \leq G$. Then for any input $x$
whose initial representation $z = f_{\theta_0}(x) \in \mathrm{col}(A)$,
the drift after $N$ steps satisfies:
\begin{equation}
    \|f_{\theta_N}(x) - f_{\theta_0}(x)\| \leq
    \underbrace{0}_{\text{first-order}} +
    \frac{L \eta^2 G^2 N(N-1)}{2}
    \label{eq:second-order-bound}
\end{equation}
\end{proposition}

\begin{proof}
By Taylor expansion of the encoder around parameters $\theta_n$:
\[
    f_{\theta_{n+1}}(x) - f_{\theta_n}(x)
    = J_{\theta_n} f(x) \cdot (-\eta g'_n) + R_n
\]
where $\|R_n\| \leq \frac{L}{2} \|\eta g'_n\|^2 \leq \frac{L\eta^2 G^2}{2}$.

Proposition~\ref{prop:zero-forgetting} shows the first-order term
$J_{\theta_n} f(x) \cdot (-\eta g'_n)$ has zero component in
$\mathrm{col}(A)$ at step $n=0$. At subsequent steps the archived
representation may have drifted by the accumulated remainder, giving
a first-order contribution of $O(\eta^2)$ per step. Summing over $N$
steps:
\[
    \|f_{\theta_N}(x) - f_{\theta_0}(x)\|
    \leq \sum_{n=0}^{N-1} \|R_n\| + O(\eta^2 N)
    \leq \frac{L \eta^2 G^2 N(N-1)}{2}
\]
\end{proof}

\begin{corollary}[Small Learning Rate Regime]
\label{cor:small-lr}
As $\eta \to 0$ with fixed total parameter displacement $\eta N = C$
(i.e.\ more steps at smaller rate), the forgetting bound becomes
$O(\eta) \to 0$. In the continuous gradient-flow limit ($\eta \to 0,
N \to \infty, \eta N = C$), the LatentShift method achieves exactly
zero forgetting for archived representations.
\end{corollary}

This explains the empirical gap between the theoretical zero-forgetting
guarantee and the small but non-zero forgetting observed on CIFAR benchmarks:
large learning rates and deep nonlinear encoders introduce non-negligible
second-order terms. The gap shrinks with smaller learning rates or
shallower encoders (as confirmed on Split-MNIST, where forgetting is
$< 1\%$).

\subsection{Connection to Hilbert's Hotel}

The LatentShift method draws an analogy with Hilbert's Grand Hotel paradox:

\begin{itemize}
    \item \textbf{Rooms} correspond to latent dimensions $\{1, 2, \ldots, d\}$.
    \item \textbf{Guests} correspond to task representations (subspace directions).
    \item The \textbf{shift operator} $n \to 2n$ in Hilbert's Hotel corresponds to
          our SVD-based subspace identification followed by orthogonal projection,
          which ``shifts'' old task representations into a protected archive and
          frees orthogonal directions for new guests.
\end{itemize}

Unlike Hilbert's Hotel (which has countably infinite rooms), the latent space
has finite dimension $d$, leading to the capacity bound in Proposition~\ref{prop:capacity}.
This boundedness motivates the lossy compression mechanism
(Proposition~\ref{prop:compression}) as a practical extension.

\textbf{Key distinction from GPM:} While GPM~\cite{saha2021gradient} also
uses orthogonal projection to prevent forgetting, it operates \emph{per-layer},
maintaining separate subspace bases for each network layer. LatentShift operates
in the \emph{holistic latent representation space}, which has several advantages:
(1)~a single, interpretable capacity measure (archive rank vs.\ latent dimension);
(2)~architectural flexibility (the projector is agnostic to encoder structure);
(3)~a cleaner theoretical framework (Propositions~\ref{prop:zero-forgetting}--\ref{prop:compression}
    apply to the unified latent space rather than fragmented per-layer spaces).
 from the main paper file.

\section{Theoretical Analysis}
\label{sec:theory}

We present formal guarantees for the LatentShift method. All proofs follow
from standard linear algebra and the properties of orthogonal projections.

\subsection{Setup and Notation}

Let $f_\theta: \mathcal{X} \to \mathbb{R}^d$ denote the encoder mapping
inputs to a $d$-dimensional latent space. After training on task $t$, we
collect latent activations $Z_t = \{f_\theta(x) : x \in \mathcal{D}_t\}$
and compute via SVD a rank-$r_t$ orthonormal basis $V_t \in \mathbb{R}^{d \times r_t}$
capturing at least $\tau$ fraction of the variance (where $\tau$ is the
threshold hyperparameter, typically 0.99).

The \emph{archive basis} after $T$ tasks is the orthonormalized union
$A_T \in \mathbb{R}^{d \times k_T}$, where $k_T \leq \sum_{t=1}^T r_t$,
obtained by QR decomposition of $[A_{T-1} \mid V_T]$.

The \emph{free-subspace projector} is:
\begin{equation}
    P = I_d - A_T A_T^\top
    \label{eq:projector}
\end{equation}

During training on task $T+1$, all gradients $g$ flowing through the latent
representation are projected as $g' = Pg$.

\subsection{Proposition 1: Zero-Forgetting Guarantee}

\begin{proposition}[Zero Forgetting]
\label{prop:zero-forgetting}
Let $A \in \mathbb{R}^{d \times k}$ be the archive basis with orthonormal
columns ($A^\top A = I_k$), and let $P = I_d - AA^\top$ be the free-subspace
projector. For any archived representation $z \in \mathrm{col}(A)$ and any
projected gradient $g' = Pg$:
\begin{equation}
    z^\top g' = 0
\end{equation}
That is, the projected gradient has zero component along any direction in
the archive subspace. Consequently, parameter updates driven by $g'$
cannot modify the encoder's output for any input whose latent
representation lies in $\mathrm{col}(A)$.
\end{proposition}

\begin{proof}
Since $z \in \mathrm{col}(A)$, there exists $\alpha \in \mathbb{R}^k$ such
that $z = A\alpha$. Then:
\begin{align}
    z^\top g' &= (A\alpha)^\top (I - AA^\top) g \\
              &= \alpha^\top A^\top g - \alpha^\top \underbrace{A^\top A}_{= I_k} A^\top g \\
              &= \alpha^\top A^\top g - \alpha^\top A^\top g \\
              &= 0
\end{align}
The first-order Taylor expansion of the encoder output change is
$\Delta f_\theta(x) \approx J_z \cdot g'$, where $J_z$ is the Jacobian.
Since $g'$ lies entirely in $\mathrm{null}(A^\top)$, and the archived
representations span $\mathrm{col}(A)$, the inner product
$z^\top \Delta f_\theta(x) = 0$ to first order.
\end{proof}

\begin{remark}[First-Order Approximation]
\label{rem:first-order}
Proposition~\ref{prop:zero-forgetting} guarantees \emph{exact} orthogonality
$z^\top g' = 0$ at the gradient level, but the claim that archived
representations are unchanged relies on a first-order Taylor expansion of
the encoder $f_\theta$. When $f_\theta$ is a deep nonlinear network,
higher-order terms in the parameter update $\Delta \theta$ can cause
small drifts in the latent representation $z = f_\theta(x)$ even though
the projected gradient $g'$ is orthogonal to the archive subspace.
We quantify this residual in Proposition~\ref{prop:second-order}.
\end{remark}

\begin{remark}[Latent-to-Parameter Gap]
\label{rem:latent-to-parameter}
Proposition~\ref{prop:zero-forgetting} establishes orthogonality in the
latent space: $z^\top g' = 0$. To connect this to parameter space, note
that the encoder $f_\theta$ maps parameters $\theta$ to latent
representations $z = f_\theta(x)$. Let $J_\theta = \partial f_\theta(x) /
\partial \theta$ denote the Jacobian of the encoder with respect to
parameters (note: this is distinct from $J_z$ used in the proof above,
which informally denotes the Jacobian of the encoder output change with
respect to latent-space perturbations). If $g' = Pg$ is the projected
gradient in latent space, then by the chain rule, the effective parameter
gradient becomes:
\begin{equation}
    \frac{\partial \mathcal{L}}{\partial \theta}\bigg|_{\text{projected}}
    = J_\theta^\top \, g'
    = J_\theta^\top \, P \, g
    \label{eq:latent-to-param}
\end{equation}
Since $g' = Pg \in \mathrm{null}(A^\top)$ by construction, the parameter
update $\Delta \theta = -\eta \, J_\theta^\top g'$ inherits the
orthogonality constraint: for any archived representation
$z = A\alpha \in \mathrm{col}(A)$, the first-order change in $z$ is
$\Delta z \approx J_\theta \, \Delta \theta = -\eta \, J_\theta J_\theta^\top g'$.
The component of $\Delta z$ along the archive is
$A^\top \Delta z = -\eta \, A^\top J_\theta J_\theta^\top g'$, which
vanishes to first order when $g' \in \mathrm{null}(A^\top)$ and
$J_\theta J_\theta^\top$ preserves this null-space structure (as holds
exactly when $J_\theta$ commutes with $P$, and approximately otherwise).
This is a \emph{first-order} guarantee; the second-order remainder is
captured by Proposition~\ref{prop:second-order}, which bounds the
cumulative drift from nonlinear effects.
\end{remark}

\subsection{Proposition 2: Capacity Bound}

\begin{proposition}[Capacity Bound]
\label{prop:capacity}
For a latent space of dimension $d$, with task $t$ occupying a subspace
of rank $r_t$, the LatentShift method guarantees zero forgetting for all
tasks if and only if:
\begin{equation}
    k_T = \mathrm{rank}(A_T) \leq d
\end{equation}
After QR reorthogonalization, $k_T \leq \sum_{t=1}^T r_t$, with equality
when task subspaces are mutually orthogonal. The maximum number of tasks
with guaranteed zero forgetting is:
\begin{equation}
    T_{\max} = \left\lfloor \frac{d}{\bar{r}} \right\rfloor
\end{equation}
where $\bar{r} = \frac{1}{T}\sum_{t=1}^T r_t$ is the average per-task rank.
When $k_T = d$, the free subspace is empty ($P = 0$) and no further
learning is possible without interference.
\end{proposition}

\begin{proof}
The archive basis $A_T$ has $k_T$ orthonormal columns in $\mathbb{R}^d$.
The free projector $P = I - A_T A_T^\top$ has rank $d - k_T$. For $P$
to be non-trivial (allowing new learning), we need $d - k_T > 0$.

Each task $t$ contributes at most $r_t$ new linearly independent directions.
The QR step ensures orthonormality and removes redundancy, so
$k_T \leq \min(\sum_t r_t, d)$. The bound is tight when task subspaces
are orthogonal.

Setting $k_T = d$ yields $P = 0$, meaning $g' = 0$ for all gradients,
preventing any parameter updates. Thus $T_{\max} = \lfloor d/\bar{r} \rfloor$.
\end{proof}

\subsection{Proposition 3: Isometry Preservation}

\begin{proposition}[Free-Subspace Isometry]
\label{prop:isometry}
The projector $P = I - AA^\top$ preserves inner products within the free
subspace. That is, for any $u, v \in \mathrm{null}(A^\top)$:
\begin{equation}
    \langle Pu, Pv \rangle = \langle u, v \rangle
\end{equation}
Moreover, $P$ preserves norms: $\|Pu\| = \|u\|$ for $u \in \mathrm{null}(A^\top)$.
\end{proposition}

\begin{proof}
If $u \in \mathrm{null}(A^\top)$, then $A^\top u = 0$, so
$Pu = u - AA^\top u = u - A \cdot 0 = u$. Similarly $Pv = v$. Therefore:
\[
    \langle Pu, Pv \rangle = \langle u, v \rangle.
\]
For norms: $\|Pu\|^2 = \langle Pu, Pu \rangle = \langle u, u \rangle = \|u\|^2$.
\end{proof}

This ensures that LatentShift does not distort the geometry of new task
representations — learning in the free subspace proceeds as if the
archive dimensions simply did not exist.

\subsection{Proposition 4: Lossy Compression Bound}

\begin{proposition}[Compression Error Bound]
\label{prop:compression}
Suppose the archive has rank $k$ with basis $A = [a_1 \mid \cdots \mid a_k]$
ordered by decreasing importance (singular values $\sigma_1 \geq \cdots
\geq \sigma_k$ from the cumulative SVD). Compressing to rank $k' < k$ by
retaining only the first $k'$ columns yields a representation error bounded by:
\begin{equation}
    \max_{z \in \mathrm{col}(A), \|z\|=1} \|z - \hat{z}\| \leq
    \sqrt{\sum_{i=k'+1}^{k} \sigma_i^2 \Big/ \sum_{i=1}^{k} \sigma_i^2}
\end{equation}
where $\hat{z}$ is the projection of $z$ onto the compressed basis
$A' = [a_1 \mid \cdots \mid a_{k'}]$.
\end{proposition}

\begin{proof}
For any unit vector $z \in \mathrm{col}(A)$, write $z = A\alpha$ with
$\|\alpha\| = 1$ (since $A$ has orthonormal columns). The projection onto
$A'$ gives $\hat{z} = A' A'^\top z = A' \alpha_{1:k'}$ where $\alpha_{1:k'}$
are the first $k'$ components. The error is:
\[
    \|z - \hat{z}\|^2 = \|\alpha_{k'+1:k}\|^2 \leq 1
\]
When the archive was built from SVD with singular values $\sigma_i$, the
importance of each direction is proportional to $\sigma_i^2$. Restricting
to the subspace captured by the cumulative SVD, the relative error for
the worst-case archived representation is bounded by the fraction of
discarded variance.
\end{proof}

\subsection{Proposition 5: Second-Order Forgetting Bound}

\begin{proposition}[Cumulative Second-Order Forgetting]
\label{prop:second-order}
Let $f_\theta$ be an $L$-smooth encoder (i.e.\ its Jacobian $J_\theta f$
is Lipschitz with constant $L$). Suppose training on task $T+1$ takes
$N$ gradient steps with learning rate $\eta$, each producing a projected
gradient $g'_n = P g_n$ with $\|g_n\| \leq G$. Then for any input $x$
whose initial representation $z = f_{\theta_0}(x) \in \mathrm{col}(A)$,
the drift after $N$ steps satisfies:
\begin{equation}
    \|f_{\theta_N}(x) - f_{\theta_0}(x)\| \leq
    \underbrace{0}_{\text{first-order}} +
    \frac{L \eta^2 G^2 N(N-1)}{2}
    \label{eq:second-order-bound}
\end{equation}
\end{proposition}

\begin{proof}
By Taylor expansion of the encoder around parameters $\theta_n$:
\[
    f_{\theta_{n+1}}(x) - f_{\theta_n}(x)
    = J_{\theta_n} f(x) \cdot (-\eta g'_n) + R_n
\]
where $\|R_n\| \leq \frac{L}{2} \|\eta g'_n\|^2 \leq \frac{L\eta^2 G^2}{2}$.

Proposition~\ref{prop:zero-forgetting} shows the first-order term
$J_{\theta_n} f(x) \cdot (-\eta g'_n)$ has zero component in
$\mathrm{col}(A)$ at step $n=0$. At subsequent steps the archived
representation may have drifted by the accumulated remainder, giving
a first-order contribution of $O(\eta^2)$ per step. Summing over $N$
steps:
\[
    \|f_{\theta_N}(x) - f_{\theta_0}(x)\|
    \leq \sum_{n=0}^{N-1} \|R_n\| + O(\eta^2 N)
    \leq \frac{L \eta^2 G^2 N(N-1)}{2}
\]
\end{proof}

\begin{corollary}[Small Learning Rate Regime]
\label{cor:small-lr}
As $\eta \to 0$ with fixed total parameter displacement $\eta N = C$
(i.e.\ more steps at smaller rate), the forgetting bound becomes
$O(\eta) \to 0$. In the continuous gradient-flow limit ($\eta \to 0,
N \to \infty, \eta N = C$), the LatentShift method achieves exactly
zero forgetting for archived representations.
\end{corollary}

This explains the empirical gap between the theoretical zero-forgetting
guarantee and the small but non-zero forgetting observed on CIFAR benchmarks:
large learning rates and deep nonlinear encoders introduce non-negligible
second-order terms. The gap shrinks with smaller learning rates or
shallower encoders (as confirmed on Split-MNIST, where forgetting is
$< 1\%$).

\subsection{Connection to Hilbert's Hotel}

The LatentShift method draws an analogy with Hilbert's Grand Hotel paradox:

\begin{itemize}
    \item \textbf{Rooms} correspond to latent dimensions $\{1, 2, \ldots, d\}$.
    \item \textbf{Guests} correspond to task representations (subspace directions).
    \item The \textbf{shift operator} $n \to 2n$ in Hilbert's Hotel corresponds to
          our SVD-based subspace identification followed by orthogonal projection,
          which ``shifts'' old task representations into a protected archive and
          frees orthogonal directions for new guests.
\end{itemize}

Unlike Hilbert's Hotel (which has countably infinite rooms), the latent space
has finite dimension $d$, leading to the capacity bound in Proposition~\ref{prop:capacity}.
This boundedness motivates the lossy compression mechanism
(Proposition~\ref{prop:compression}) as a practical extension.

\textbf{Key distinction from GPM:} While GPM~\cite{saha2021gradient} also
uses orthogonal projection to prevent forgetting, it operates \emph{per-layer},
maintaining separate subspace bases for each network layer. LatentShift operates
in the \emph{holistic latent representation space}, which has several advantages:
(1)~a single, interpretable capacity measure (archive rank vs.\ latent dimension);
(2)~architectural flexibility (the projector is agnostic to encoder structure);
(3)~a cleaner theoretical framework (Propositions~\ref{prop:zero-forgetting}--\ref{prop:compression}
    apply to the unified latent space rather than fragmented per-layer spaces).
 from the main paper file.

\section{Theoretical Analysis}
\label{sec:theory}

We present formal guarantees for the LatentShift method. All proofs follow
from standard linear algebra and the properties of orthogonal projections.

\subsection{Setup and Notation}

Let $f_\theta: \mathcal{X} \to \mathbb{R}^d$ denote the encoder mapping
inputs to a $d$-dimensional latent space. After training on task $t$, we
collect latent activations $Z_t = \{f_\theta(x) : x \in \mathcal{D}_t\}$
and compute via SVD a rank-$r_t$ orthonormal basis $V_t \in \mathbb{R}^{d \times r_t}$
capturing at least $\tau$ fraction of the variance (where $\tau$ is the
threshold hyperparameter, typically 0.99).

The \emph{archive basis} after $T$ tasks is the orthonormalized union
$A_T \in \mathbb{R}^{d \times k_T}$, where $k_T \leq \sum_{t=1}^T r_t$,
obtained by QR decomposition of $[A_{T-1} \mid V_T]$.

The \emph{free-subspace projector} is:
\begin{equation}
    P = I_d - A_T A_T^\top
    \label{eq:projector}
\end{equation}

During training on task $T+1$, all gradients $g$ flowing through the latent
representation are projected as $g' = Pg$.

\subsection{Proposition 1: Zero-Forgetting Guarantee}

\begin{proposition}[Zero Forgetting]
\label{prop:zero-forgetting}
Let $A \in \mathbb{R}^{d \times k}$ be the archive basis with orthonormal
columns ($A^\top A = I_k$), and let $P = I_d - AA^\top$ be the free-subspace
projector. For any archived representation $z \in \mathrm{col}(A)$ and any
projected gradient $g' = Pg$:
\begin{equation}
    z^\top g' = 0
\end{equation}
That is, the projected gradient has zero component along any direction in
the archive subspace. Consequently, parameter updates driven by $g'$
cannot modify the encoder's output for any input whose latent
representation lies in $\mathrm{col}(A)$.
\end{proposition}

\begin{proof}
Since $z \in \mathrm{col}(A)$, there exists $\alpha \in \mathbb{R}^k$ such
that $z = A\alpha$. Then:
\begin{align}
    z^\top g' &= (A\alpha)^\top (I - AA^\top) g \\
              &= \alpha^\top A^\top g - \alpha^\top \underbrace{A^\top A}_{= I_k} A^\top g \\
              &= \alpha^\top A^\top g - \alpha^\top A^\top g \\
              &= 0
\end{align}
The first-order Taylor expansion of the encoder output change is
$\Delta f_\theta(x) \approx J_z \cdot g'$, where $J_z$ is the Jacobian.
Since $g'$ lies entirely in $\mathrm{null}(A^\top)$, and the archived
representations span $\mathrm{col}(A)$, the inner product
$z^\top \Delta f_\theta(x) = 0$ to first order.
\end{proof}

\begin{remark}[First-Order Approximation]
\label{rem:first-order}
Proposition~\ref{prop:zero-forgetting} guarantees \emph{exact} orthogonality
$z^\top g' = 0$ at the gradient level, but the claim that archived
representations are unchanged relies on a first-order Taylor expansion of
the encoder $f_\theta$. When $f_\theta$ is a deep nonlinear network,
higher-order terms in the parameter update $\Delta \theta$ can cause
small drifts in the latent representation $z = f_\theta(x)$ even though
the projected gradient $g'$ is orthogonal to the archive subspace.
We quantify this residual in Proposition~\ref{prop:second-order}.
\end{remark}

\begin{remark}[Latent-to-Parameter Gap]
\label{rem:latent-to-parameter}
Proposition~\ref{prop:zero-forgetting} establishes orthogonality in the
latent space: $z^\top g' = 0$. To connect this to parameter space, note
that the encoder $f_\theta$ maps parameters $\theta$ to latent
representations $z = f_\theta(x)$. Let $J_\theta = \partial f_\theta(x) /
\partial \theta$ denote the Jacobian of the encoder with respect to
parameters (note: this is distinct from $J_z$ used in the proof above,
which informally denotes the Jacobian of the encoder output change with
respect to latent-space perturbations). If $g' = Pg$ is the projected
gradient in latent space, then by the chain rule, the effective parameter
gradient becomes:
\begin{equation}
    \frac{\partial \mathcal{L}}{\partial \theta}\bigg|_{\text{projected}}
    = J_\theta^\top \, g'
    = J_\theta^\top \, P \, g
    \label{eq:latent-to-param}
\end{equation}
Since $g' = Pg \in \mathrm{null}(A^\top)$ by construction, the parameter
update $\Delta \theta = -\eta \, J_\theta^\top g'$ inherits the
orthogonality constraint: for any archived representation
$z = A\alpha \in \mathrm{col}(A)$, the first-order change in $z$ is
$\Delta z \approx J_\theta \, \Delta \theta = -\eta \, J_\theta J_\theta^\top g'$.
The component of $\Delta z$ along the archive is
$A^\top \Delta z = -\eta \, A^\top J_\theta J_\theta^\top g'$, which
vanishes to first order when $g' \in \mathrm{null}(A^\top)$ and
$J_\theta J_\theta^\top$ preserves this null-space structure (as holds
exactly when $J_\theta$ commutes with $P$, and approximately otherwise).
This is a \emph{first-order} guarantee; the second-order remainder is
captured by Proposition~\ref{prop:second-order}, which bounds the
cumulative drift from nonlinear effects.
\end{remark}

\subsection{Proposition 2: Capacity Bound}

\begin{proposition}[Capacity Bound]
\label{prop:capacity}
For a latent space of dimension $d$, with task $t$ occupying a subspace
of rank $r_t$, the LatentShift method guarantees zero forgetting for all
tasks if and only if:
\begin{equation}
    k_T = \mathrm{rank}(A_T) \leq d
\end{equation}
After QR reorthogonalization, $k_T \leq \sum_{t=1}^T r_t$, with equality
when task subspaces are mutually orthogonal. The maximum number of tasks
with guaranteed zero forgetting is:
\begin{equation}
    T_{\max} = \left\lfloor \frac{d}{\bar{r}} \right\rfloor
\end{equation}
where $\bar{r} = \frac{1}{T}\sum_{t=1}^T r_t$ is the average per-task rank.
When $k_T = d$, the free subspace is empty ($P = 0$) and no further
learning is possible without interference.
\end{proposition}

\begin{proof}
The archive basis $A_T$ has $k_T$ orthonormal columns in $\mathbb{R}^d$.
The free projector $P = I - A_T A_T^\top$ has rank $d - k_T$. For $P$
to be non-trivial (allowing new learning), we need $d - k_T > 0$.

Each task $t$ contributes at most $r_t$ new linearly independent directions.
The QR step ensures orthonormality and removes redundancy, so
$k_T \leq \min(\sum_t r_t, d)$. The bound is tight when task subspaces
are orthogonal.

Setting $k_T = d$ yields $P = 0$, meaning $g' = 0$ for all gradients,
preventing any parameter updates. Thus $T_{\max} = \lfloor d/\bar{r} \rfloor$.
\end{proof}

\subsection{Proposition 3: Isometry Preservation}

\begin{proposition}[Free-Subspace Isometry]
\label{prop:isometry}
The projector $P = I - AA^\top$ preserves inner products within the free
subspace. That is, for any $u, v \in \mathrm{null}(A^\top)$:
\begin{equation}
    \langle Pu, Pv \rangle = \langle u, v \rangle
\end{equation}
Moreover, $P$ preserves norms: $\|Pu\| = \|u\|$ for $u \in \mathrm{null}(A^\top)$.
\end{proposition}

\begin{proof}
If $u \in \mathrm{null}(A^\top)$, then $A^\top u = 0$, so
$Pu = u - AA^\top u = u - A \cdot 0 = u$. Similarly $Pv = v$. Therefore:
\[
    \langle Pu, Pv \rangle = \langle u, v \rangle.
\]
For norms: $\|Pu\|^2 = \langle Pu, Pu \rangle = \langle u, u \rangle = \|u\|^2$.
\end{proof}

This ensures that LatentShift does not distort the geometry of new task
representations — learning in the free subspace proceeds as if the
archive dimensions simply did not exist.

\subsection{Proposition 4: Lossy Compression Bound}

\begin{proposition}[Compression Error Bound]
\label{prop:compression}
Suppose the archive has rank $k$ with basis $A = [a_1 \mid \cdots \mid a_k]$
ordered by decreasing importance (singular values $\sigma_1 \geq \cdots
\geq \sigma_k$ from the cumulative SVD). Compressing to rank $k' < k$ by
retaining only the first $k'$ columns yields a representation error bounded by:
\begin{equation}
    \max_{z \in \mathrm{col}(A), \|z\|=1} \|z - \hat{z}\| \leq
    \sqrt{\sum_{i=k'+1}^{k} \sigma_i^2 \Big/ \sum_{i=1}^{k} \sigma_i^2}
\end{equation}
where $\hat{z}$ is the projection of $z$ onto the compressed basis
$A' = [a_1 \mid \cdots \mid a_{k'}]$.
\end{proposition}

\begin{proof}
For any unit vector $z \in \mathrm{col}(A)$, write $z = A\alpha$ with
$\|\alpha\| = 1$ (since $A$ has orthonormal columns). The projection onto
$A'$ gives $\hat{z} = A' A'^\top z = A' \alpha_{1:k'}$ where $\alpha_{1:k'}$
are the first $k'$ components. The error is:
\[
    \|z - \hat{z}\|^2 = \|\alpha_{k'+1:k}\|^2 \leq 1
\]
When the archive was built from SVD with singular values $\sigma_i$, the
importance of each direction is proportional to $\sigma_i^2$. Restricting
to the subspace captured by the cumulative SVD, the relative error for
the worst-case archived representation is bounded by the fraction of
discarded variance.
\end{proof}

\subsection{Proposition 5: Second-Order Forgetting Bound}

\begin{proposition}[Cumulative Second-Order Forgetting]
\label{prop:second-order}
Let $f_\theta$ be an $L$-smooth encoder (i.e.\ its Jacobian $J_\theta f$
is Lipschitz with constant $L$). Suppose training on task $T+1$ takes
$N$ gradient steps with learning rate $\eta$, each producing a projected
gradient $g'_n = P g_n$ with $\|g_n\| \leq G$. Then for any input $x$
whose initial representation $z = f_{\theta_0}(x) \in \mathrm{col}(A)$,
the drift after $N$ steps satisfies:
\begin{equation}
    \|f_{\theta_N}(x) - f_{\theta_0}(x)\| \leq
    \underbrace{0}_{\text{first-order}} +
    \frac{L \eta^2 G^2 N(N-1)}{2}
    \label{eq:second-order-bound}
\end{equation}
\end{proposition}

\begin{proof}
By Taylor expansion of the encoder around parameters $\theta_n$:
\[
    f_{\theta_{n+1}}(x) - f_{\theta_n}(x)
    = J_{\theta_n} f(x) \cdot (-\eta g'_n) + R_n
\]
where $\|R_n\| \leq \frac{L}{2} \|\eta g'_n\|^2 \leq \frac{L\eta^2 G^2}{2}$.

Proposition~\ref{prop:zero-forgetting} shows the first-order term
$J_{\theta_n} f(x) \cdot (-\eta g'_n)$ has zero component in
$\mathrm{col}(A)$ at step $n=0$. At subsequent steps the archived
representation may have drifted by the accumulated remainder, giving
a first-order contribution of $O(\eta^2)$ per step. Summing over $N$
steps:
\[
    \|f_{\theta_N}(x) - f_{\theta_0}(x)\|
    \leq \sum_{n=0}^{N-1} \|R_n\| + O(\eta^2 N)
    \leq \frac{L \eta^2 G^2 N(N-1)}{2}
\]
\end{proof}

\begin{corollary}[Small Learning Rate Regime]
\label{cor:small-lr}
As $\eta \to 0$ with fixed total parameter displacement $\eta N = C$
(i.e.\ more steps at smaller rate), the forgetting bound becomes
$O(\eta) \to 0$. In the continuous gradient-flow limit ($\eta \to 0,
N \to \infty, \eta N = C$), the LatentShift method achieves exactly
zero forgetting for archived representations.
\end{corollary}

This explains the empirical gap between the theoretical zero-forgetting
guarantee and the small but non-zero forgetting observed on CIFAR benchmarks:
large learning rates and deep nonlinear encoders introduce non-negligible
second-order terms. The gap shrinks with smaller learning rates or
shallower encoders (as confirmed on Split-MNIST, where forgetting is
$< 1\%$).

\subsection{Connection to Hilbert's Hotel}

The LatentShift method draws an analogy with Hilbert's Grand Hotel paradox:

\begin{itemize}
    \item \textbf{Rooms} correspond to latent dimensions $\{1, 2, \ldots, d\}$.
    \item \textbf{Guests} correspond to task representations (subspace directions).
    \item The \textbf{shift operator} $n \to 2n$ in Hilbert's Hotel corresponds to
          our SVD-based subspace identification followed by orthogonal projection,
          which ``shifts'' old task representations into a protected archive and
          frees orthogonal directions for new guests.
\end{itemize}

Unlike Hilbert's Hotel (which has countably infinite rooms), the latent space
has finite dimension $d$, leading to the capacity bound in Proposition~\ref{prop:capacity}.
This boundedness motivates the lossy compression mechanism
(Proposition~\ref{prop:compression}) as a practical extension.

\textbf{Key distinction from GPM:} While GPM~\cite{saha2021gradient} also
uses orthogonal projection to prevent forgetting, it operates \emph{per-layer},
maintaining separate subspace bases for each network layer. LatentShift operates
in the \emph{holistic latent representation space}, which has several advantages:
(1)~a single, interpretable capacity measure (archive rank vs.\ latent dimension);
(2)~architectural flexibility (the projector is agnostic to encoder structure);
(3)~a cleaner theoretical framework (Propositions~\ref{prop:zero-forgetting}--\ref{prop:compression}
    apply to the unified latent space rather than fragmented per-layer spaces).
